\providecommand{\Needspace}[1]{}
\Needspace{8\baselineskip}
\item \textbf{Caída de un paquete soltado desde un dron con viento horizontal}.
\nopagebreak[4]

\begin{tcolorbox}[colback=blue!2!white,colframe=blue!55!black,title={Enunciado},breakable]
Un dron vuela horizontalmente y suelta un paquete desde una altura inicial $h_0$. Mientras el paquete cae, el viento produce una aceleración horizontal constante $a_w$. Se desea estudiar el movimiento usando arreglos de NumPy y visualizar, en una misma figura, la posición, velocidad y aceleración en las direcciones horizontal y vertical.

Considere los siguientes valores constantes:
\[
h_0 = 80 \ \text{m},
\qquad
v_{x0} = 12 \ \text{m/s},
\qquad
v_{y0} = 0 \ \text{m/s},
\qquad
a_w = 0.8 \ \text{m/s}^2,
\qquad
g = 9.8 \ \text{m/s}^2.
\]

El tiempo debe representarse mediante un arreglo temporal construido con \texttt{np.linspace}. En este problema use un arreglo \texttt{t} con $300$ valores entre $t=0$ y $t=5$ segundos.

Las ecuaciones cinemáticas del movimiento son:
\[
x(t) = v_{x0}t + \frac{1}{2}a_wt^2,
\]
\[
y(t) = h_0 + v_{y0}t - \frac{1}{2}gt^2.
\]

Las velocidades en cada dirección están dadas por:
\[
v_x(t) = v_{x0} + a_wt,
\]
\[
v_y(t) = v_{y0} - gt.
\]

Las aceleraciones correspondientes son:
\[
a_x(t) = a_w,
\]
\[
a_y(t) = -g.
\]

Escriba un programa en Python que realice lo siguiente:

\begin{itemize}
    \item[(a)] Defina las constantes físicas del problema y construya el arreglo temporal $t$ usando \texttt{np.linspace} con inicio $0$, final $5$ y $300$ puntos.

    \item[(b)] Calcule los arreglos de posición $x(t)$ y $y(t)$ usando las ecuaciones cinemáticas entregadas.

    \item[(c)] Calcule los arreglos de velocidad $v_x(t)$ y $v_y(t)$.

    \item[(d)] Encuentre el índice y el valor mínimo de la posición vertical $y(t)$. Además, calcule el promedio de los arreglos de velocidad $v_x(t)$ y $v_y(t)$.

    \item[(e)] Calcule los arreglos de aceleración $a_x(t)$ y $a_y(t)$. Ambas aceleraciones deben representarse como arreglos con la misma dimensión que $t$ para poder graficarlas correctamente.

    \item[(f)] Realice una figura con subgráficos usando \texttt{plt.subplots(3, 2)}. En la primera columna debe graficar las cantidades asociadas al movimiento vertical:
    \[
    y(t) \ \text{vs.} \ t,
    \qquad
    v_y(t) \ \text{vs.} \ t,
    \qquad
    a_y(t) \ \text{vs.} \ t.
    \]
    En la segunda columna debe graficar las cantidades asociadas al movimiento horizontal:
    \[
    x(t) \ \text{vs.} \ t,
    \qquad
    v_x(t) \ \text{vs.} \ t,
    \qquad
    a_x(t) \ \text{vs.} \ t.
    \]

    \item[(g)] Agregue títulos, nombres de ejes y una presentación limpia a cada gráfico, de modo que la figura permita comparar claramente la evolución horizontal y vertical del paquete.
\end{itemize}
\end{tcolorbox}

\begin{solution}
Una posible solución consiste en calcular todas las cantidades como arreglos de NumPy. Las aceleraciones constantes también se escriben como arreglos para que puedan graficarse junto con el tiempo.

\begin{pythonjupyter}
import numpy as np
import matplotlib.pyplot as plt

# Constantes físicas
h0 = 80       # altura inicial en metros
vx0 = 12      # velocidad horizontal inicial en m/s
vy0 = 0       # velocidad vertical inicial en m/s
aw = 0.8      # aceleración horizontal por viento en m/s^2
g = 9.8       # aceleración de gravedad en m/s^2

# Arreglo temporal
t = np.linspace(0, 5, 300)

# Posiciones
x = vx0 * t + 0.5 * aw * t**2
y = h0 + vy0 * t - 0.5 * g * t**2

# Velocidades
vx = vx0 + aw * t
vy = vy0 - g * t

# Índice y valor mínimo de y
idx_y_min = np.argmin(y)
y_min = y[idx_y_min]

# Promedios de velocidad
vx_prom = np.mean(vx)
vy_prom = np.mean(vy)

print("Posición vertical mínima:")
print("Índice:", idx_y_min)
print("Valor:", y_min)

print("\nPromedios:")
print("Promedio de vx:", vx_prom)
print("Promedio de vy:", vy_prom)

# Aceleraciones
ax = aw * np.ones_like(t)
ay = -g * np.ones_like(t)

# Figura con subgráficos
fig, axs = plt.subplots(3, 2, figsize=(10, 9))

# Columna izquierda: movimiento vertical
axs[0, 0].plot(t, y)
axs[0, 0].set_title("Posición vertical")
axs[0, 0].set_xlabel("t [s]")
axs[0, 0].set_ylabel("y(t) [m]")
axs[0, 0].grid(True)

axs[1, 0].plot(t, vy)
axs[1, 0].set_title("Velocidad vertical")
axs[1, 0].set_xlabel("t [s]")
axs[1, 0].set_ylabel("vy(t) [m/s]")
axs[1, 0].grid(True)

axs[2, 0].plot(t, ay)
axs[2, 0].set_title("Aceleración vertical")
axs[2, 0].set_xlabel("t [s]")
axs[2, 0].set_ylabel("ay(t) [m/s^2]")
axs[2, 0].grid(True)

# Columna derecha: movimiento horizontal
axs[0, 1].plot(t, x)
axs[0, 1].set_title("Posición horizontal")
axs[0, 1].set_xlabel("t [s]")
axs[0, 1].set_ylabel("x(t) [m]")
axs[0, 1].grid(True)

axs[1, 1].plot(t, vx)
axs[1, 1].set_title("Velocidad horizontal")
axs[1, 1].set_xlabel("t [s]")
axs[1, 1].set_ylabel("vx(t) [m/s]")
axs[1, 1].grid(True)

axs[2, 1].plot(t, ax)
axs[2, 1].set_title("Aceleración horizontal")
axs[2, 1].set_xlabel("t [s]")
axs[2, 1].set_ylabel("ax(t) [m/s^2]")
axs[2, 1].grid(True)

plt.tight_layout()
plt.show()
\end{pythonjupyter}
\end{solution}
